\chapter{Практика}

\section{Задача}
Проанализировать пакеты DNS в Wireshark.

\section{Анализ с помощью Wireshark}

Выполним 

\begin{lstlisting}
    sudo nslookup yandex.ru
\end{lstlisting}

А в Wireshark будем ловить пакеты (в фильтре можно установить dns).

\subsection{Запрос}

\begin{figure}[H] 
    \centering
    \includegraphics[width=0.8\textwidth]{whr.png}
    \caption{Запрос к серверу}
\end{figure}

Важные поля:

\begin{enumerate}
    \item Questions - число запросов
    \item Queries - запись, к которой нужно найти адрес.
\end{enumerate}

\subsection{Ответ}

\begin{figure}[H] 
    \centering
    \includegraphics[width=0.8\textwidth]{wh.png}
    \caption{Ответ от сервера}
\end{figure}

Важные поля:

\begin{enumerate}
    \item Answer RRs - число основных записей в ответе.
    \item Authority RRs - число авторитетных записей.
    \item Additional RRs - число доп. записей.
    \item Answers, Authoritative nameservers, Additional records - здесь перечислены записи, которые вернул сервер. Самая важная информация в каждой записи - IP.
\end{enumerate}

\endinput